\documentclass[12pt]{article}
\usepackage{graphicx} % Required for inserting images
\usepackage[margin=0.5in]{geometry}
\usepackage{amsmath, amssymb}
\usepackage{mathrsfs}


\usepackage{soul}


% Dr. David Brown Commands
\newcommand{\ds}{\displaystyle}
\newcommand{\R}{\mathbb{R}}
\newcommand{\N}{\mathbb{N}}
\newcommand{\Q}{\mathbb{Q}}
\newcommand{\C}{\mathbb{C}}


% Commands to get script r
\usepackage{calligra}
\DeclareMathAlphabet{\mathcalligra}{T1}{calligra}{m}{n}
\DeclareFontShape{T1}{calligra}{m}{n}{<->s*[2.2]callig15}{}
\newcommand{\scripty}[1]{\ensuremath{\mathcalligra{#1}}}

% Unit commands
\newcommand{\degree}{\text{°}}
\newcommand{\unitSpeed}{\frac{\text{m}}{\text{s}}}
\newcommand{\unitAcceleration}{\frac{\text{m}}{\text{s}^2}}

% Constant commands
\newcommand{\avogadro}{6.022\times10^{23}}

\newcommand{\boltzmannConstK}{1.381\times10^{-23}\frac{\text{J}}{\text{K}}}
\newcommand{\planckConst}{6.626\times10^{-34}\text{J}\cdot\text{s}}

\newcommand{\atmP}{1.01325\times10^5\,\text{Pa}}

% Commands to easily write partial derivatives
\newcommand{\partDeriv}[2]{\frac{\partial #1}{\partial #2}}
\newcommand{\doublePartDeriv}[2]{\frac{\partial^2 #1}{\partial^2 #2}}


\title{Problem 45}
\author{Gabriel Decker}


\begin{document}



\maketitle
\begin{flushleft}



Your apparent weight is less than your actual weight when you feel less force on you that normally counteracts the force of gravity.
For example, in an elevator going up, you feel heavier when it begins accelerating upwards, and you feel lighter when it begins deaccelerating when it's near the top.
So if you feel your apparent weight to be 70\% of the normal weight, then the plane must be accelerating downwards, putting less force on you.
Say you have a mass $m$. When the plane is moving normally (not varying in the y-direction), the net force in the y-direction is zero.
The force of gravity $\vec{F}_g$ and the normal force $\vec{F}_n$ cancel out, so
$$\vec{F}_g+\vec{F}_n=m\vec{a}=0$$
Since $\vec{F}_g=-mg\hat{y}$, this implies that $\vec{F}_n=mg\hat{y}$.\\~\\

However, in this case, you feel 70\% of your apparent weight, so this implies that 30\% of the regular normal force is missing, so $\vec{F}_n=\frac{7}{10}mg\hat{y}$.
Note that now, the forces don't cancel out, so you (and the plane) must be accelerating.
We can use Newton's 2nd again to find how much we're accelerating in the y-direction.
$$\vec{F}_g+\vec{F}_n=m\vec{a}$$
$$\implies-mg\hat{y}+\frac{7}{10}mg\hat{y}=m\vec{a}$$
$$\implies\vec{a}=-\frac{3}{10}g\hat{y}$$\\~\\

The negative shows that we're accelerating downward.
Since we're dealing with one component, the magnitude is simply
$$a=\frac{3}{10}g=2.94\,\text{m/s}^2$$




\end{flushleft}

\end{document}
